\chapter{Unital versus non-unital rings}

\fontsize{12}{14.5}\selectfont

In this chapter, we aim to introduce a motivating example for the notion of ideally exact category. This is based on an article by George Janelidze \cite{Janelidze2024}.

Consider the category $\mathsf{Ring}$ of rings with unit. The notion of ideal is, at first glance, somewhat unnatural in the context of categories. In fact, in general, an ideal of a ring with unit is not an object in the category $\mathsf{Ring}$: if an ideal contains the unit, then it coincides with the whole ring. However, we observe that it is an object of the category $\mathsf{Ring_*}$ of not necessarily unital rings.

It is known that $\mathsf{Ring}$ and $\mathsf{Ring_*}$ are exact and protomodular categories (see \cite{BorceuxBourn2004} for example). They both have binary coproducts, even if they are not the same. The coproduct in the category $\mathsf{Ring_*}$ is the free product, while in $\mathsf{Ring}$ it is the free product with amalgamation, in order to ensure that there is a unit in the coproduct.

That said, they have a crucial difference: $\mathsf{Ring}$ is not pointed, whereas $\mathsf{Ring_*}$ is. In fact, in $\mathsf{Ring_*}$ the initial and terminal objects are both the zero ring $\{0\}$, while in $\mathsf{Ring}$ the initial object $0$ is the ring $\mathbb{Z}$ of integers, whereas the terminal object $1$ is the trivial ring $\{0\}$. Thus, we conclude the $\mathsf{Ring_*}$ is semiabelian, while $\mathsf{Ring}$ fails to be.

There is an obvious inclusion functor (that forgets the existence of a unit element):
\[U: \mathsf{Ring} \to \mathsf{Ring_*}\]
This functor has a left adjoint: $\mathbb{Z} \ltimes (-)$. This functor, is given by
\begin{align*}
F: \mathsf{Ring_*} &\to \mathsf{Ring} \\
R &\mapsto \mathbb{Z} \ltimes R
\end{align*}
where $\mathbb{Z} \ltimes R$ is $\mathbb{Z} \times R$ as a set, with componentwise addition, while the product $\cdot$ is defined as follows
\begin{align*}
\cdot : (\mathbb{Z} \ltimes R) \times (\mathbb{Z} \ltimes R) &\to \mathbb{Z} \ltimes R \\
((m,r),(n,s))&\mapsto (mn, nr + ms + rs)
\end{align*}
The unit of the ring $(\mathbb{Z} \ltimes R)$ is $(1,0_R)$. This construction gives us a unital ring.

For a morphism $f: R \to S$ in $\mathsf{Ring_*}$,
\begin{align*}
    F(f) : \mathbb{Z} \ltimes R &\to \mathbb{Z} \ltimes S \\
    (n,r) &\mapsto (n, f(r))
\end{align*}

The unit and the counit of the adjunction are given by
\begin{align*}
\eta_R: R &\to UF(R)=\mathbb{Z} \ltimes R \\
r &\mapsto (0,r)
\end{align*}
\begin{align*}
\epsilon_R: FU(R)=\mathbb{Z} \ltimes R &\to R  \\
(n,r) &\mapsto (n1_R + r)
\end{align*}
One can check that if $f: R \to
U(S)$ is a morphism in $\mathsf{Ring_*}$, with $S$ unital, then the function 
\begin{align*}
Z \ltimes R &\to S \\
(n,r) &\mapsto (n1_S + f(r))
\end{align*}
gives a morphism in $\mathsf{Ring}$, so that the unit $\eta$ has the universal property required for an adjunction.

The functor $U$ turns out to be monadic, simply by Beck's theorem, since coequalizers in $\mathsf{Ring}$ and $\mathsf{Ring_*}$ are the same. In fact, given $f,g$ morphisms, the coequalizer of $(f,g)$ is the canonical projection
\[\pi: B \to B/I \ \ \ \ \ \ \ I = (f(r) - g(r) \ | \ r \in R)\]
Notice that the unit $\eta$ of the adjunction is cartesian, i.e. the naturality squares are pullbacks. In fact, one can easily prove that, given a morphism $f:R \to R'$ in $\mathsf{Ring_*}$, the pullback in $\mathsf{Ring_*}$ of $\eta_{R'}$ and $UF(f)$
\begin{align*}
R'&\times_{UF(R')} UF(R) = R'\times_{\mathbb{Z}\ltimes R'} (\mathbb{Z}\ltimes R) = \\[1.5ex]
&= \{ (r',(n,r)) \in R'\times (\mathbb{Z}\ltimes R) \ | \  UF(f)(n,r) = \eta_{R'}(r') \} \\[1.5ex]
&=\{ (r',(n,r)) \in R'\times (\mathbb{Z}\ltimes R) \ | \ n=0 \ \land \ f(r)=r' \} \\[1.5ex]
&=\{ (f(r),(0,r)) \ | \  r \in R \}
\end{align*}
is isomorphic to $R$.
\[
% https://tikzcd.yichuanshen.de/#N4Igdg9gJgpgziAXAbVABwnAlgFyxMJZABgBpiBdUkANwEMAbAVxiRAAIAlAcgB1e8AW3gB9YP0F0cACwBGs4AC0AvvwhpmcLt2XsAFBKlyFKtRqZbOAShDLS6TLnyEUARlIAmKrUYs2PW3sQDGw8AiIAZlJXb3pmVkQQQxl5JVVedU0uQIdQ50jPWN8EpN5JFJN0zIttHODHMJdkdxjqOL9EzltvGCgAc3giUAAzACcIQSQyEBwIJHcQWRgwKCQAWgjp9pLk6VHBYDRxgCtlEVcQagY6JYYABQb8xNGsPukcOrGJqepZpA9qEsVkhNm1imxdvtDiczh5LiBrrcHnlwokGDBhh87CNxpNENM-ogACxXLBgEpQOhwaS9T64pAkmZzRALbZsYZ074s37MqI+eIQ3gwHB0MQ8ZScvGMwkA-kdUrC0VdbEgL542WEvlsxIAVQAYnphjZlBRlEA
\begin{tikzcd}
 R'\times_{\mathbb{Z}\ltimes R'} (\mathbb{Z}\ltimes R) \arrow[rdd, "\mathrm{proj}_1"', bend right] \arrow[rrrd, "\mathrm{proj}_2", bend left] \arrow[rd, dashed] &                                       &  &                                       \\
                                                                                                                                                               & R \arrow[d, "f"] \arrow[rr, "\eta_R"] &  & \mathbb{Z}\ltimes R \arrow[d, "UF(f)"] \\
                                                                                                                                                               & R' \arrow[rr, "\eta_{R'}"]            &  & \mathbb{Z}\ltimes R'                  
\end{tikzcd}
\]










From the properties of the adjunction above, Janelidze extrapolated the
notion of ideally exact category. In order to better understand the link
between the guiding example and the abstract notion, we need a few more
observations.

First of all, we have an obvious equivalence $(\mathsf{Ring} \downarrow 1)\cong \mathsf{Ring}$, as for
any category with a terminal object. More interestingly, there is another equivalence of categories

\[
(\mathsf{Ring} \downarrow 0)\cong (\mathsf{Ring} \downarrow \mathbb{Z}) \cong \mathsf{Ring_*}
\]
given by the unitalization functor 
\begin{align*}
    (F,\pi_\mathbb{Z}) : \mathsf{Ring_*} &\to (\mathsf{Ring} \downarrow \mathbb{Z})\\
    R &\mapsto (\mathbb{Z} \ltimes R, \pi_\mathbb{Z}:\mathbb{Z} \ltimes R \to \mathbb{Z})
\end{align*}
and its quasi inverse $\mathrm{Ker}$
\begin{align*}
    \mathrm{Ker} : (\mathsf{Ring} \downarrow \mathbb{Z}) &\to \mathsf{Ring_*}\\
    (R,\phi:R \to \mathbb{Z}) &\mapsto \mathrm{Ker}(\phi) 
\end{align*}
Indeed, for an $R$ in $\mathsf{Ring_*}$, $\mathrm{Ker}(\pi_\mathbb{Z} : \mathbb{Z} \ltimes R \to \mathbb{Z})$
is clearly isomorphic to $R$.

On the other hand, given a $\phi: R \to \mathbb{Z}$ in $\mathsf{Ring}$, 
and fixed a kernel $K$ of $\phi$, the functions
\begin{align*}
    \mathbb{Z} \ltimes K &\to R \\
    (n,k) &\mapsto n1_R + k
\end{align*}
\begin{align*}
    R &\to \mathbb{Z} \ltimes K \\
    r &\mapsto (\phi(r), r - \phi(r)1_R)
\end{align*}
are inverse each other and provide an isomorphism
$(\mathbb{Z} \ltimes K, \pi_\mathbb{Z}) \cong (R,\phi)$
in $(\mathsf{Ring} \downarrow \mathbb{Z})$.

Notice that $\mathrm{Ker}$ is used here in the usual algebraic sense,
and is not a kernel in $\mathsf{Ring}$ in the categorical sense, since the category is not
pointed, and more importantly, if $\phi \neq 0$ then $\mathrm{Ker}(\phi)$ is not a unital subring of $R$. More appropriately, one should regard it as $\mathrm{Ker}(U(\phi))$, so that it becomes
a categorical kernel in the
(pointed) category $\mathsf{Ring_*}$. This perspective will be
essential to define ideals in the ideally exact context.

Based on the discussion above, one may regard $U$ as a functor from $(\mathsf{Ring} \downarrow 1)$ to $(\mathsf{Ring} \downarrow 0)$, up to equivalence. Actually, we can be even more precise.
For each $R$ in $\mathsf{Ring}$, the functions
\begin{align*}
\mathbb{Z} \times R &\to \mathbb{Z} \ltimes R \\
(n,r) &\mapsto (n, -n1_R + r)
\end{align*}
\begin{align*}
\mathbb{Z} \ltimes R &\to \mathbb{Z} \times R \\
(n,r) &\mapsto (n, n1_R + r)
\end{align*}
are inverse each other and provide an isomorphism $(\mathbb{Z} \times K, \pi_\mathbb{Z}) \cong (\mathbb{Z} \ltimes K, \pi_\mathbb{Z})$ in $(\mathsf{Ring} \downarrow \mathbb{Z})$. But $\pi_\mathbb{Z}: \mathbb{Z} \times K \to \mathbb{Z}$ is also the first projection of the pullback

\[
% https://tikzcd.yichuanshen.de/#N4Igdg9gJgpgziAXAbVABwnAlgFyxMJZABgBpiBdUkANwEMAbAVxiRGIB0O8BbeAfQCMAAgBKIAL6l0mXPkIoATKUFVajFm0GTpIDNjwEiZVdXrNWidjpkH5RZZTMbLIcRLUwoAc3hFQAGYAThA8SADM1DgQSILOFmwAhAB67rrBoUjKINGx8ZpWiTYgGWGIZDkxiJHqCVZcPHQ4ABZBPMBoIQBWEvyKxaVIFbmI2eYFIA1Nre2dED1CkhQSQA
\begin{tikzcd}
\mathbb{Z} \times R \arrow[rr, "\mathrm{proj}_2"] \arrow[d, "\mathrm{proj}_1"] &  & R \arrow[d, "!^R"] \\
0 = \mathbb{Z} \arrow[rr, "!"]                                                      &  & 1 = \{0\}                
\end{tikzcd}
\]
in $\mathsf{Ring}$. In other words, the composite functor
\[
(\mathsf{Ring} \downarrow 1) \xrightarrow{\quad\cong\quad} \mathsf{Ring} \xrightarrow{\quad U\quad} \mathsf{Ring_*} \xrightarrow{(F,\pi_\mathbb{Z})} (\mathsf{Ring} \downarrow 0)
\]
is isomorphic to the pullback functor \Cref{chap:pb-funct}.
\[
!^* : (\mathsf{Ring} \downarrow 1) \to (\mathsf{Ring} \downarrow 0)
\]
To sum up, there is an equivalence of adjunctions:
\[
% https://tikzcd.yichuanshen.de/#N4Igdg9gJgpgziAXAbVABwnAlgFyxMJZABgBpiBdUkANwEMAbAVxiRAAoAdTgWzpwAWAIwBmwAEpYwAcwC+AAm5QIAdzB0AThtXyAjAEoQs0uky58hFGQBMVWoxZt5XXv2FjJMhUtXqtO7j5BISFgAC1ZQ2NTbDwCImtyO3pmVkQQQLdRCSk5IxMQDFiLBNJbahTHdMzgj1yAfQAqWSM7GChpeCJQEW0eJDIQHAgkXWo4ASwRHCQAWjH7VLYAQgA9RvyevtHqYYHxyem5hcq0kGX65c2QXoh+xEShkcQAZgOpmcR5ioczgFVrrd7m8nkhHhMPscfkt0gAxQHbRCDPYPaFVDKcADGBGkIF2dCwDDYkDArGoAhgdCgSDATAYDGoDCkZzgECZ1MZdCEMAYAAUzHFLCAGDAjtEboiFiiQRCjogTr82NxsTI8cKuTz+cV4ukRWKKLIgA
\begin{tikzcd}
(\mathsf{Ring} \downarrow 1) \arrow[dd, "!^*", shift left=3] \arrow[rr, leftrightarrow, "\cong"] \arrow[dd, phantom, "\dashv" description]                       &  & \mathsf{Ring} \arrow[dd, "U", shift left=3] \arrow[dd, phantom, "\dashv" description]      \\
                                                                                                      &  &                                             \\
 (\mathsf{Ring} \downarrow 0) \arrow[uu, "!_!", shift left=3] \arrow[rr, "\cong", leftrightarrow] &  & \mathsf{Ring_*} \arrow[uu, "F", shift left=3]
\end{tikzcd}
\]

It is precisely this interpretation of the adjunction $F \dashv U$ in terms of
pullback functor that suggests how to extrapolate from the example of rings
the more general notion of ideally exact category, illustrated in \Cref{chap:ideally-ex}.